\section{Unbounded complexes}
\label{section-unbounded}

\noindent
A reference for the material in this section is \cite{Spaltenstein}.
The following lemma is useful to find ``good'' left resolutions of
unbounded complexes.

\begin{lemma}
\label{lemma-special-direct-system}
Let $\mathcal{A}$ be an abelian category. Let
$\mathcal{P} \subset \Ob(\mathcal{A})$ be a subset.
Assume $\mathcal{P}$ contains $0$, is closed under (finite) direct sums,
and every object of $\mathcal{A}$ is a quotient of an
element of $\mathcal{P}$. Let $K^\bullet$ be a complex.
There exists a commutative diagram
$$
\xymatrix{
P_1^\bullet \ar[d] \ar[r] & P_2^\bullet \ar[d] \ar[r] & \ldots \\
\tau_{\leq 1}K^\bullet \ar[r] & \tau_{\leq 2}K^\bullet \ar[r] & \ldots
}
$$
in the category of complexes such that
\begin{enumerate}
\item the vertical arrows are quasi-isomorphisms and termwise surjective,
\item $P_n^\bullet$ is a bounded above complex with terms in
$\mathcal{P}$,
\item the arrows $P_n^\bullet \to P_{n + 1}^\bullet$
are termwise split injections and each cokernel
$P^i_{n + 1}/P^i_n$ is an element of $\mathcal{P}$.
\end{enumerate}
\end{lemma}

\begin{proof}
We are going to use that the homotopy category $K(\mathcal{A})$ is a
triangulated category, see Proposition
\ref{proposition-homotopy-category-triangulated}.
By Lemma \ref{lemma-subcategory-left-resolution} we can find a
termwise surjective map of complexes $P_1^\bullet \to \tau_{\leq 1}K^\bullet$
which is a quasi-isomorphism
such that the terms of $P_1^\bullet$ are in $\mathcal{P}$.
By induction it suffices, given
$P_1^\bullet, \ldots, P_n^\bullet$ to construct
$P_{n + 1}^\bullet$ and the maps
$P_n^\bullet \to P_{n + 1}^\bullet$ and
$P_{n + 1}^\bullet \to \tau_{\leq n + 1}K^\bullet$.

\medskip\noindent
Choose a distinguished triangle
$P_n^\bullet \to \tau_{\leq n + 1}K^\bullet \to C^\bullet \to P_n^\bullet[1]$
in $K(\mathcal{A})$. Applying
Lemma \ref{lemma-subcategory-left-resolution} we choose a
map of complexes $Q^\bullet \to C^\bullet$
which is a quasi-isomorphism such that the terms of $Q^\bullet$
are in $\mathcal{P}$. By the axioms of triangulated categories
we may fit the composition $Q^\bullet \to C^\bullet \to P_n^\bullet[1]$ into
a distinguished triangle
$P_n^\bullet \to P_{n + 1}^\bullet \to Q^\bullet \to P_n^\bullet[1]$
in $K(\mathcal{A})$.
By Lemma \ref{lemma-improve-distinguished-triangle-homotopy}
we may and do assume
$0 \to P_n^\bullet \to P_{n + 1}^\bullet \to Q^\bullet \to 0$
is a termwise split short exact sequence. This implies that
the terms of $P_{n + 1}^\bullet$ are in $\mathcal{P}$ and that
$P_n^\bullet \to P_{n + 1}^\bullet$ is a termwise split injection
whose cokernels are in $\mathcal{P}$.
By the axioms of triangulated categories we obtain a map
of distinguished triangles
$$
\xymatrix{
P_n^\bullet \ar[r] \ar[d] &
P_{n + 1}^\bullet \ar[r] \ar[d] &
Q^\bullet \ar[r] \ar[d] &
P_n^\bullet[1] \ar[d] \\
P_n^\bullet \ar[r] &
\tau_{\leq n + 1}K^\bullet \ar[r] &
C^\bullet \ar[r] &
P_n^\bullet[1]
}
$$
in the triangulated category $K(\mathcal{A})$. Choose an actual morphism of
complexes $f : P_{n + 1}^\bullet \to \tau_{\leq n + 1}K^\bullet$.
The left square of the diagram above commutes up to homotopy, but as
$P_n^\bullet \to P_{n + 1}^\bullet$ is a termwise split injection
we can lift the homotopy and modify our choice of $f$ to make it commute.
Finally, $f$ is a quasi-isomorphism, because both $P_n^\bullet \to P_n^\bullet$
and $Q^\bullet \to C^\bullet$ are.

\medskip\noindent
At this point we have all the properties we want, except we don't know
that the map $f : P_{n + 1}^\bullet \to \tau_{\leq n + 1}K^\bullet$
is termwise surjective. Since we have the commutative diagram
$$
\xymatrix{
P_n^\bullet \ar[d] \ar[r] & P_{n + 1}^\bullet \ar[d] \\
\tau_{\leq n}K^\bullet \ar[r] & \tau_{\leq n + 1}K^\bullet
}
$$
of complexes, by induction hypothesis we see that $f$ is surjective
on terms in all degrees except possibly $n$ and $n + 1$. Choose
an object $P \in \mathcal{P}$ and a surjection $q : P \to K^n$.
Consider the map
$$
g :
P^\bullet = (\ldots \to 0 \to P \xrightarrow{1} P \to 0 \to \ldots)
\longrightarrow
\tau_{\leq n + 1}K^\bullet
$$
with first copy of $P$ in degree $n$ and maps given by
$q$ in degree $n$ and $d_K \circ q$ in degree $n + 1$.
This is a surjection in degree $n$ and the cokernel in
degree $n + 1$ is $H^{n + 1}(\tau_{\leq n + 1}K^\bullet)$;
to see this recall that $\tau_{\leq n + 1}K^\bullet$ has
$\Ker(d_K^{n + 1})$ in degree $n + 1$.
However, since $f$ is a quasi-isomorphism we know that
$H^{n + 1}(f)$ is surjective. Hence after replacing 
$f : P_{n + 1}^\bullet \to \tau_{\leq n + 1}K^\bullet$
by
$f \oplus g : P_{n + 1}^\bullet \oplus P^\bullet \to \tau_{\leq n + 1}K^\bullet$
we win.
\end{proof}

\noindent
In some cases we can use the lemma above to show that a left derived
functor is everywhere defined.

\begin{proposition}
\label{proposition-left-derived-exists}
Let $F : \mathcal{A} \to \mathcal{B}$ be a right exact functor
of abelian categories. Let $\mathcal{P} \subset \Ob(\mathcal{A})$ be a
subset. Assume
\begin{enumerate}
\item $\mathcal{P}$ contains $0$, is closed under (finite) direct sums,
and every object of $\mathcal{A}$ is a quotient of an
element of $\mathcal{P}$,
\item for any bounded above acyclic complex $P^\bullet$ of
$\mathcal{A}$ with $P^n \in \mathcal{P}$ for all $n$ the
complex $F(P^\bullet)$ is exact,
\item $\mathcal{A}$ and $\mathcal{B}$ have colimits
of systems over $\mathbf{N}$,
\item colimits over $\mathbf{N}$ are exact in both
$\mathcal{A}$ and $\mathcal{B}$, and
\item $F$ commutes with colimits over $\mathbf{N}$.
\end{enumerate}
Then $LF$ is defined on all of $D(\mathcal{A})$.
\end{proposition}

\begin{proof}
By (1) and Lemma \ref{lemma-subcategory-left-resolution} for any bounded
above complex $K^\bullet$ there exists a quasi-isomorphism
$P^\bullet \to K^\bullet$ with $P^\bullet$ bounded above and
$P^n \in \mathcal{P}$ for all $n$. Suppose that
$s : P^\bullet \to (P')^\bullet$ is a quasi-isomorphism of bounded
above complexes consisting of objects of $\mathcal{P}$. Then
$F(P^\bullet) \to F((P')^\bullet)$ is a quasi-isomorphism because
$F(C(s)^\bullet)$ is acyclic by assumption (2). This already shows that
$LF$ is defined on $D^{-}(\mathcal{A})$ and that a bounded above
complex consisting of objects of $\mathcal{P}$ computes $LF$, see
Lemma \ref{lemma-find-existence-computes}.

\medskip\noindent
Next, let $K^\bullet$ be an arbitrary complex of $\mathcal{A}$.
Choose a diagram
$$
\xymatrix{
P_1^\bullet \ar[d] \ar[r] & P_2^\bullet \ar[d] \ar[r] & \ldots \\
\tau_{\leq 1}K^\bullet \ar[r] & \tau_{\leq 2}K^\bullet \ar[r] & \ldots
}
$$
as in Lemma \ref{lemma-special-direct-system}. Note that
the map $\colim P_n^\bullet \to K^\bullet$ is a quasi-isomorphism
because colimits over $\mathbf{N}$ in $\mathcal{A}$ are exact
and $H^i(P_n^\bullet) = H^i(K^\bullet)$ for $n > i$. We claim that
$$
F(\colim P_n^\bullet) = \colim F(P_n^\bullet)
$$
(termwise colimits) is $LF(K^\bullet)$, i.e., that $\colim P_n^\bullet$
computes $LF$. To see this, by Lemma \ref{lemma-find-existence-computes},
it suffices to prove the following claim. Suppose that
$$
\colim Q_n^\bullet = Q^\bullet
\xrightarrow{\ \alpha\ }
P^\bullet = \colim P_n^\bullet
$$
is a quasi-isomorphism of complexes, such that each
$P_n^\bullet$, $Q_n^\bullet$ is a bounded above complex whose terms are
in $\mathcal{P}$ and the maps $P_n^\bullet \to \tau_{\leq n}P^\bullet$ and
$Q_n^\bullet \to \tau_{\leq n}Q^\bullet$ are quasi-isomorphisms.
Claim: $F(\alpha)$ is a quasi-isomorphism.

\medskip\noindent
The problem is that we do not assume that $\alpha$ is given as a colimit
of maps between the complexes $P_n^\bullet$ and $Q_n^\bullet$. However,
for each $n$ we know that the solid arrows in the diagram
$$
\xymatrix{
& R^\bullet \ar@{..>}[d] \\
P_n^\bullet \ar[d] &
L^\bullet \ar@{..>}[l] \ar@{..>}[r] &
Q_n^\bullet \ar[d] \\
\tau_{\leq n}P^\bullet \ar[rr]^{\tau_{\leq n}\alpha} & &
\tau_{\leq n}Q^\bullet
}
$$
are quasi-isomorphisms. Because quasi-isomorphisms form a multiplicative
system in $K(\mathcal{A})$ (see Lemma \ref{lemma-acyclic})
we can find a quasi-isomorphism
$L^\bullet \to P_n^\bullet$ and map of complexes $L^\bullet \to Q_n^\bullet$
such that the diagram above commutes up to homotopy. Then
$\tau_{\leq n}L^\bullet \to L^\bullet$ is a quasi-isomorphism.
Hence (by the first part of the proof) we can find a bounded above
complex $R^\bullet$ whose terms are in $\mathcal{P}$ and a quasi-isomorphism
$R^\bullet \to L^\bullet$ (as indicated in the diagram). Using the result
of the first paragraph of the proof we see that
$F(R^\bullet) \to F(P_n^\bullet)$ and $F(R^\bullet) \to F(Q_n^\bullet)$
are quasi-isomorphisms. Thus we obtain an isomorphism
$H^i(F(P_n^\bullet)) \to H^i(F(Q_n^\bullet))$ fitting into the commutative
diagram
$$
\xymatrix{
H^i(F(P_n^\bullet)) \ar[r] \ar[d] &
H^i(F(Q_n^\bullet)) \ar[d] \\
H^i(F(P^\bullet)) \ar[r] &
H^i(F(Q^\bullet))
}
$$
The exact same argument shows that these maps are also compatible
as $n$ varies. Since by (4) and (5) we have
$$
H^i(F(P^\bullet)) =
H^i(F(\colim P_n^\bullet)) =
H^i(\colim F(P_n^\bullet)) = \colim H^i(F(P_n^\bullet))
$$
and similarly for $Q^\bullet$ we conclude that
$H^i(\alpha) : H^i(F(P^\bullet) \to H^i(F(Q^\bullet)$ is an isomorphism
and the claim follows.
\end{proof}

\begin{lemma}
\label{lemma-special-inverse-system}
Let $\mathcal{A}$ be an abelian category. Let
$\mathcal{I} \subset \Ob(\mathcal{A})$ be a subset.
Assume $\mathcal{I}$ contains $0$, is closed under (finite) products,
and every object of $\mathcal{A}$ is a subobject of an
element of $\mathcal{I}$. Let $K^\bullet$ be a complex.
There exists a commutative diagram
$$
\xymatrix{
\ldots \ar[r] &
\tau_{\geq -2}K^\bullet \ar[r] \ar[d] &
\tau_{\geq -1}K^\bullet \ar[d] \\
\ldots \ar[r] & I_2^\bullet \ar[r] & I_1^\bullet
}
$$
in the category of complexes such that
\begin{enumerate}
\item the vertical arrows are quasi-isomorphisms and termwise injective,
\item $I_n^\bullet$ is a bounded below complex with terms in $\mathcal{I}$,
\item the arrows $I_{n + 1}^\bullet \to I_n^\bullet$ are termwise split
surjections and $\Ker(I^i_{n + 1} \to I^i_n)$ is an element of $\mathcal{I}$.
\end{enumerate}
\end{lemma}

\begin{proof}
This lemma is dual to
Lemma \ref{lemma-special-direct-system}.
\end{proof}







